% !PRINTED 397
\weberpage{397}
%<<CONT
\[
Q_{n-1} \leqq Q_n,
\]
wo das Gleichheitszeichen nur im Falle des unteren Zeichens,
also im Falle der Gleichung (9) möglich ist.

Ferner ist (da $b < \sqrt{D}$ ist)
\[
Q_{n-1} = \frac{t-ub}{2} > \frac{t-u\sqrt{D}}{2} = \frac{t^2-u^2D}{2(t+u\sqrt{D})} = \frac{\mp 2}{t+u\sqrt{D}},
\]
also $Q_{n-1}$ Null oder positiv, und Null kann nur im Falle des
oberen Zeichens, also im Falle der Gleichung (8) auftreten.
Daraus schliessen wir
\weq{12}{0 \leqq Q_{n-1} \leqq Q_n,}
wo das Gleichheitszeichen in der unteren Grenze nur im Falle
der Gleichung (8), also der oberen Zeichen, in der oberen Grenze
nur im Falle der Gleichung (9), also der unteren Zeichen vorkommen kann.

Wir entwickeln nun den rationalen Bruch $P_n : Q_n$ in einen
endlichen Kettenbruch
\[
\frac{P_n}{Q_n} = (\alpha_0, \alpha_1, \alpha_2 \ldots \alpha_{n-1}),
\]
indem wir $n$ nach \S.\,115, I, so annehmen, dass
\[
(-1)^n = \mp 1
\]
wird. Bedeutet dann $P' : Q'$ den vorletzten Näherungsbruch, so ist
\[
P_n Q' - Q_n P' = \mp 1
\]
und es ist nach \S.\,117, II:
\[
0 \leqq Q' \leqq Q_n,
\]
wo das Gleichheitszeichen in der unteren Grenze nur für $n=1$,
also im Falle der oberen Zeichen, in der oberen Grenze nur für
$n=2$, also der unteren Zeichen vorkommen kann, ebenso wie
in (12). Daraus folgt nach \S.\,118, III:
\[
Q' = Q_{n-1}, \quad P' = P_{n-1},
\]
und es ist also nach (10)
\[
(\alpha_0, \alpha_1, \alpha_2 \ldots \alpha_{n-1}, \omega) = \omega,
\]
d. h. $[\alpha_0, \alpha_1, \alpha_2 \ldots \alpha_{n-1}]$ ist eine Periode des Kettenbruches für
$\omega$, aus der nach den Formeln (4), (5) die Lösung $t, u$ von (7)
hergeleitet wird. Damit ist bewiesen, dass wir auf diese Weise
\gesperrt{alle Lösungen} der Gleichung (7) finden.

Hieraus lassen sich noch einige wichtige Folgerungen ziehen.
